\documentclass[10pt, letterpaper]{article}

% Inhaltsverzeichnis für Pakettypen (nur für Übersicht im Header, wird nicht im Dokument angezeigt)
% 1. Seitenlayout und Ränder
% 2. Sprache und Zeichensatz
% 3. Mathematik und Theorem-Umgebungen
% 4. Eigene Makros
% 5. Diagramme und Grafiken
% 6. Tabellen und Aufzählungen
% 7. Inhaltsverzeichnis
% 8. Abschnittsüberschriften
% 9. Abstrakt-Umgebung
% 10. Todos/Notizen
% 11. Rahmen/Box-Umgebungen
% 12. Python-Integration
% 13. Literaturverwaltung
% 14. Hyperlinks
% 15. Absatzeinstellungen
% 16. Umgebungen
% 17  Graphik
% 18  Extra
% 00. Titel und Autor

\usepackage{slashed}

% --- 1. Seitenlayout und Ränder ---
\usepackage[margin=3cm]{geometry}

% --- 2. Sprache und Zeichensatz ---
\usepackage[english]{babel}
\usepackage[T1]{fontenc}
\usepackage[utf8]{inputenc}

% --- 3. Mathematik und Theorem-Umgebungen ---
\usepackage{amsmath, amssymb, amsthm}
\usepackage{mathrsfs}
\DeclareMathOperator{\WF}{WF}

% --- 4. Eigene Makros ---
\usepackage{xcolor}
\newcommand{\SKP}{\langle\cdot,\cdot\rangle}
\newcommand{\id}{\operatorname{id}}
\newcommand{\sgn}{\operatorname{sgn}}
\newcommand{\Ad}{\operatorname{Ad}}
\newcommand{\ad}{\operatorname{ad}}
\newcommand{\R}{\mathbb{R}}
\newcommand{\N}{\mathbb{N}}
\newcommand{\Q}{\mathbb{Q}}
\newcommand{\Z}{\mathbb{Z}}
\newcommand{\C}{\mathbb{C}}
\newcommand{\QU}{\mathbb{H}}
\newcommand{\Cinf}{C^{\infty}}
\newcommand{\Mod}{\mathrm{Mod}}
\newcommand{\Alg}{\mathrm{Alg}}
\newcommand{\Diff}{\mathrm{Diff}}
\newcommand{\Iso}{\mathrm{Iso}}
\newcommand{\Aut}{\mathrm{Aut}}
\newcommand{\End}{\mathrm{End}}
\newcommand{\Hom}{\mathrm{Hom}}
\newcommand{\Cl}{\mathrm{Cl}}
\newcommand{\entwurf}[1]{\textcolor{red}{#1}}
\newcommand{\GL}{\mathrm{GL}}
\newcommand{\Mat}{\mathrm{Mat}}
\newcommand{\SL}{\mathrm{SL}}
\newcommand{\SO}{\mathrm{SO}}
\newcommand{\SU}{\mathrm{SU}}
\newcommand{\Sp}{\mathrm{Sp}}
\newcommand{\Spin}{\mathrm{Spin}}
\newcommand{\Pin}{\mathrm{Pin}}
\newcommand{\Lor}{\mathrm{Lor}}
\newcommand{\Ogroup}{\mathrm{O}}
\newcommand{\U}{\mathrm{U}}

% --- 5. Diagramme und Grafiken ---
\usepackage{graphicx}
\graphicspath{{../../Library/Images/}}
\usepackage[export]{adjustbox}
\usepackage{tikz}
\usetikzlibrary{decorations.pathreplacing, arrows.meta, positioning}
\usepackage{tikz-cd}

% --- 6. Tabellen und Aufzählungen ---
\usepackage{enumitem}
\setlist[itemize]{left=0.5cm}

\newenvironment{romanenum}[1][]
  {%
    \ifx&#1&
    \else
      \textbf{#1}\quad
    \fi
    \begin{enumerate}[label=\roman*)]
  }
  {%
    \end{enumerate}%
  }

% --- 7. Inhaltsverzeichnis ---
\usepackage{tocloft}
\renewcommand{\cftsecfont}{\footnotesize}
\renewcommand{\cftsubsecfont}{\footnotesize}
\renewcommand{\cftsubsubsecfont}{\footnotesize}
\renewcommand{\cftsecpagefont}{\footnotesize}
\renewcommand{\cftsubsecpagefont}{\footnotesize}
\renewcommand{\cftsubsubsecpagefont}{\footnotesize}
\usepackage{etoc}

% --- 8. Abschnittsüberschriften ---
\usepackage{titlesec}
\titleformat{\section}{\normalfont\large\bfseries}{\thesection}{1em}{}
\titleformat{\subsection}{\normalfont\normalsize\bfseries}{\thesubsection}{0.5em}{}
\titleformat{\subsubsection}{\normalfont\normalsize\bfseries}{\thesubsubsection}{0.5em}{}
\setcounter{secnumdepth}{4}

% --- 9. Abstrakt-Umgebung ---
\usepackage{changepage}
\renewenvironment{abstract}
  {
    \begin{adjustwidth}{1.5cm}{1.5cm}
    \small
    \textsc{Abstract. –}%
  }
  {
    \end{adjustwidth}
  }

% --- 10. Todos/Notizen ---
\usepackage{todonotes}
% Hellblau definieren
\definecolor{hellblau}{rgb}{0.3,0.6,1}
\newenvironment{note}{\par\begingroup\color{hellblau}\itshape}{\par\endgroup}

% --- 11. Rahmen/Box-Umgebungen ---
\usepackage{mdframed}
\usepackage{tcolorbox}
\colorlet{shadecolor}{gray!25}

\newenvironment{customTheorem}
  {\vspace{10pt}%
   \begin{mdframed}[
     backgroundcolor=gray!20,
     linewidth=0pt,
     innertopmargin=10pt,
     innerbottommargin=10pt,
     skipabove=\dimexpr\topsep+\ht\strutbox\relax,
     skipbelow=\topsep,
   ]}
  {\end{mdframed}
   \vspace{10pt}%
  }

% --- 12. Python-Integration ---
% (Deaktiviert in dieser Version, aktiviere bei Bedarf)
% \usepackage{pythontex}
% \usepackage[makestderr]{pythontex}

% --- 13. Literaturverwaltung ---
\usepackage{csquotes}
\usepackage[backend=biber, style=alphabetic, citestyle=alphabetic]{biblatex}
\addbibresource{bibliography.bib}

% --- 14. Hyperlinks ---
\usepackage{hyperref}
\hypersetup{
  colorlinks   = true,
  urlcolor     = blue,
  linkcolor    = blue,
  citecolor    = blue,
  frenchlinks  = true
}

% --- 15. Absatzeinstellungen ---
\usepackage[parfill]{parskip}
\sloppy

% --- 16. Umgebungen ---
\usepackage{thmtools}

\newcommand{\CustomHeading}[3]{%
  \par\medskip\noindent%
  \textbf{#1 #2} \textnormal{(#3)}.\enskip%
}

\newenvironment{DEF}[2]{\begin{unitbox}\CustomHeading{Definition}{#1}{#2}}{\end{unitbox}}
\newenvironment{PROP}[2]{\begin{unitbox}\CustomHeading{Proposition}{#1}{#2}}{\end{unitbox}}
\newenvironment{THEO}[2]{\begin{unitbox}\CustomHeading{Theorem}{#1}{#2}}{\end{unitbox}}
\newenvironment{LEM}[2]{\begin{unitbox}\CustomHeading{Lemma}{#1}{#2}}{\end{unitbox}}
\newenvironment{KORO}[2]{\begin{unitbox}\CustomHeading{Corollar}{#1}{#2}}{\end{unitbox}}
\newenvironment{REM}[2]{\begin{unitbox}\CustomHeading{Remark}{#1}{#2}}{\end{unitbox}}
\newenvironment{EXA}[2]{\begin{unitbox}\CustomHeading{Example}{#1}{#2}}{\end{unitbox}}
\newenvironment{STUD}[2]{\begin{unitbox}\CustomHeading{Study}{#1}{#2}}{\end{unitbox}}
\newenvironment{CONC}[2]{\begin{unitbox}\CustomHeading{Concept}{#1}{#2}}{\end{unitbox}}
\newenvironment{OTH}[2]{\begin{unitbox}\CustomHeading{Other}{#1}{#2}}{\end{unitbox}}
\newenvironment{EXE}[2]{\begin{unitbox}\CustomHeading{Exercise}{#1}{#2}}{\end{unitbox}}
\newenvironment{MOT}[2]{\begin{unitbox}\CustomHeading{Motivation}{#1}{#2}}{\end{unitbox}}
\newenvironment{PROOF}[2]{\begin{unitbox}\CustomHeading{Proof}{#1}{#2}}{\end{unitbox}}

% --- Unit Umgebung für Source-Inhalte ---
\usepackage{mdframed}
\newmdenv[
  linewidth=1pt,
  topline=false,
  bottomline=false,
  rightline=false,
  leftmargin=0cm,
  rightmargin=0cm,
  skipabove=10pt,
  skipbelow=10pt,
  innertopmargin=0.5\baselineskip,
  innerbottommargin=0.5\baselineskip,
  backgroundcolor=gray!10,
  linecolor=gray
]{unitbox}

\newenvironment{unit}[1]
  {\begin{unitbox}\textbf{Unit #1}\par\smallskip}
  {\end{unitbox}}

% --- 17. ... ---

% --- 18. Extras ---
\usepackage{stmaryrd}
\usepackage{bbold}  % falls du \mathbb{1} nutzen willst

% --- 00. Titel und Autor ---
\title{Mein Titel}
\author{Tim Jaschik}
\date{\today}

\begin{document}

\maketitle
\rule{\textwidth}{0.5pt}
\begin{abstract}
Kurze Beschreibung …
\end{abstract}
\rule{\textwidth}{0.5pt}
\vspace{0.5cm}

\tableofcontents

\pagebreak

\subsection{Formelle Charakterisierung der Speziellen Relativitätstheorie (SRT)}

\emph{Motivation.}
Die klassische Mechanik beschreibt physikalische Prozesse in der euklidischen Raumzeit
$\R^3\times\R$ mit absoluter Zeit.  
Experimentell (Michelson–Morley, Maxwell-Theorie) zeigte sich jedoch, 
dass die Lichtgeschwindigkeit $c$ unabhängig vom Beobachter ist.  
Dies erfordert eine Raumzeitstruktur, in der
\[
\text{(i) keine ausgezeichnete Zeitrichtung existiert,}\qquad
\text{(ii) $c$ eine absolute Invariante ist.}
\]
Diese Forderungen führen zu einer geometrischen Raumzeitstruktur mit \emph{Minkowski-Metrik}.
Die SRT ist daher keine Dynamik, sondern eine \emph{Raumzeit-Theorie} —  
sie liefert den geometrischen Rahmen, in dem Dynamiken (Felder, Teilchenbewegungen) 
formuliert werden.

\bigskip
\emph{1. Geometrische Grundstruktur.}
\begin{itemize}
\item Die physikalische Arena ist die \emph{Minkowski-Raumzeit}
\[
(M,g) = (\R^{1,3}, \eta),
\quad
\eta = \mathrm{diag}(1,-1,-1,-1),
\]
eine glatte vierdimensionale Mannigfaltigkeit mit pseudo\-riemann\-scher Metrik $\eta$.
\item Die Metrik definiert die \emph{kausale Struktur}:
\[
ds^2 = \eta_{\mu\nu}dx^\mu dx^\nu
=
c^2dt^2 - d\mathbf{x}^2.
\]
\item Ereignisse $p,q\in M$ heißen
zeitartig ($ds^2>0$), lichtartig ($ds^2=0$), oder raumartig ($ds^2<0$).
\item Die Menge aller linearen Isometrien von $(M,\eta)$ bildet die \emph{Lorentzgruppe}
\[
O(1,3)
= \{\Lambda\in \mathrm{GL}(4,\R)\mid \Lambda^T\eta\Lambda = \eta\}.
\]
Deren zusammenhängende und orientierungs-/zeit-orientierte Komponente ist
\[
SO^+(1,3) = \{\Lambda\in O(1,3)\mid \det\Lambda=+1,\ \Lambda^0{}_0>0\}.
\]
\item Ergänzt durch Raumzeit-Translationen $(a,\Lambda)$ entsteht die \emph{Poincaré-Gruppe}
\[
\mathcal P = \R^{1,3}\rtimes SO^+(1,3),
\]
die vollständige Isometriegruppe der Minkowski-Raumzeit.
\end{itemize}

\bigskip
\emph{2. Physikalische Symmetrieprinzipien.}
\begin{itemize}
\item (\textbf{Relativitätsprinzip})  
Die Form physikalischer Gesetze ist in allen Inertialsystemen gleich.
\item (\textbf{Konstanz der Lichtgeschwindigkeit})  
Für alle Beobachter gilt $|v_{\text{Licht}}| = c$.
\end{itemize}
Diese beiden Prinzipien implizieren die Existenz einer Metrik $\eta$, deren Isometriegruppe
$\mathcal P$ die Symmetrie aller Naturgesetze bildet.

\bigskip
\emph{3. Formeller Kern: Natürlichkeit physikalischer Operatoren.}

Ein physikalisches Gesetz ist eine Gleichung
\[
D(\phi) = 0,
\]
wobei
\begin{itemize}
\item $\phi$ ein physikalisches Feld (z.\,B. elektromagnetisch, spinorial, skalars) auf $M$ ist,
\item $D$ ein geometrischer Differentialoperator (z.\,B. $\Box_\eta+m^2$, $i\slashed{\partial}-m$) ist.
\end{itemize}

\textbf{Forderung der Forminvarianz (Naturtreue):}
Für jede Isometrie $f\in\mathrm{Iso}(M,\eta)$ gilt
\[
\boxed{
\quad f^*\circ D = D\circ f^* \quad
\text{auf allen Feldern }\phi.
}
\]
Damit:
\[
D(\phi)=0 \quad\Longleftrightarrow\quad D(f^*\phi)=0.
\]
Das Gesetz ist also unabhängig von der Wahl des Inertialsystems.

\bigskip
\emph{4. Funktorielle Formulierung.}

Wir definieren die Kategorie
\[
\mathbf{Sys}_{\eta} :=
\{(M,\eta,\mathcal F,D)\mid (M,\eta)\ \text{Minkowski-Raumzeit},\ 
\mathcal F\to M\ \text{Feldbündel},\ D:J^k\mathcal F\to \mathcal E \text{ natürlich}\}.
\]

Ein Diffeomorphismus $f:(M,\eta)\to(M,\eta)$ mit $f^*\eta=\eta$ 
induziert Bündelabbildungen und Pullbacks $f^*:\mathfrak F\to\mathfrak F$.
Dann ist die SRT charakterisiert durch die Funktorialität:
\[
D_{f(U)}\circ J^k(f^*) = f^*\circ D_U.
\]

\textbf{Interpretation:}  
Die spezielle Relativitätstheorie ist die Forderung, dass die physikalischen 
Gesetze $D(\phi)=0$ auf der Minkowski-Raumzeit natürliche Transformationen unter 
den Isometrien $(M,\eta)$ sind.

\bigskip
\emph{5. Physikalische Bedeutung.}
\begin{itemize}
\item Die SRT stellt den \emph{universellen Raumzeitrahmen} bereit, in dem 
alle physikalischen Prozesse lokal stattfinden.  
Felder $\phi$ modellieren physikalische Größen (Materie, Strahlung, Energieverteilung).
\item Die Theorie selbst liefert keine Dynamik — sie \emph{definiert die Invarianzstruktur}, 
innerhalb derer alle Dynamiken (Maxwell, Klein–Gordon, Dirac, Yang–Mills) formuliert werden.
\item Die Metrik $\eta$ und die Poincaré-Gruppe kodieren alle experimentell bestätigten 
Symmetrien freier Raumzeitprozesse:
\begin{align*}
&\text{(Translationsinvarianz)} &&\Rightarrow \text{Impuls- und Energieerhaltung},\\
&\text{(Lorentz-Invarianz)} &&\Rightarrow \text{Erhaltung von Drehimpuls und Boostgeneratoren}.
\end{align*}
\item Somit ist die SRT kein einzelnes Gesetz, sondern die \emph{geometrische Theorie} 
aller Prozesse, deren Dynamik Lorentz-kovariant ist.
\end{itemize}

\bigskip
\emph{6. Formale Charakterisierung in einem Satz.}

Die spezielle Relativitätstheorie ist die Theorie physikalischer Systeme $(\mathcal F,D)$ auf $(\R^{1,3},\eta)$, deren Gesetze natürliche Operatoren unter der Poincaré-Gruppe sind.

Oder äquivalent:
\[
\forall f\in\mathrm{Iso}(\R^{1,3},\eta):\quad
f^*\circ D = D\circ f^*,
\qquad
\text{d.h.}\quad D(\phi)=0 \iff D(f^*\phi)=0.
\]

\bigskip
\emph{7. Physikalischer Gehalt.}
\begin{itemize}
\item \emph{Raumzeitstruktur:} Minkowski-Geometrie $(M,\eta)$, Lichtkegelstruktur invariant.
\item \emph{Symmetrien:} Poincaré-Gruppe $\mathcal P$ als Gruppe der Beobachterwechsel.
\item \emph{Physikalische Prozesse:} Geodätische, Teilchenbewegungen, Feldkonfigurationen
$\phi$ auf $(M,\eta)$.
\item \emph{Theoretischer Inhalt:} Alle physikalischen Gleichungen $D(\phi)=0$
sind Poincaré-invariant, d.\,h.:
\[
f^*\eta=\eta \ \Rightarrow\ f^*(D\phi)=D(f^*\phi).
\]
\end{itemize}

\bigskip
\textbf{In einem Satz:}

Die spezielle Relativitätstheorie ist die Geometrisierung des Relativitätsprinzips: Physik = Tensoranalysis auf $(\R^{1,3},\eta)$.

\pagebreak

\subsection{Vom experimentellen Denken zur Forminvarianz $f^*\circ D = D\circ f^*$}

\emph{Ziel.} 
Wir entwickeln eine bottom–up-Argumentationskette von konkreten, experimentell motivierten
Überlegungen hin zur abstrakten Forderung
\[
\forall f\in\mathrm{Iso}(\R^{1,3},\eta):\quad
f^*\circ D = D\circ f^*,\qquad
\text{also}\qquad D(\phi)=0 \iff D(f^*\phi)=0,
\]
ohne anfänglich Tensorrechnung vorauszusetzen. Hierbei steht $f$ für einen zulässigen
\emph{Beobachterwechsel} (Isometrie), $\phi$ für eine physikalische Größe (Feld/Konfiguration),
und $D$ für den \emph{Gesetzesoperator} (Gleichung).

\bigskip
\emph{A. Startpunkt: Was soll „gleiche Physik aus verschiedenen Blickwinkeln“ heißen?}

\textbf{A1. Beobachter und Messprozeduren.}
Ein \emph{Bezugssystem} (Inertialsystem) ist eine komplette Messanordnung aus synchronisierten Uhren
und starren Maßstäben. Zwei Beobachter $S,S'$ vergleichen Messungen desselben Vorgangs
(freier Fall, Lichtimpuls, Teilchendekay, Interferenzmuster).

\textbf{A2. Empirische Forderung 1 (Vergleichbarkeit).}
Wenn $S$ eine Versuchsanordnung beschreibt (Geometrie des Aufbaus, Startzeitpunkt, gemessene Werte),
muss $S'$ eine \emph{systematische} Vorschrift besitzen, wie er: 
(i) dieselben Ereignisse identifiziert, 
(ii) deren Messwerte berechnet.
\emph{Konsequenz:} Es braucht eine Abbildung $f$ auf der Ereignismenge (Raumzeit), die „Gleiches zu Gleichem“ zuordnet.

\textbf{A3. Empirische Forderung 2 (Lichtkonstanz \& Kausalität).}
Lichtlaufzeiten, Koinzidenzen („Detektor triggert beim Eintreffen des Photons“),
und die Einteilung in zeit-/raum-/lichtartige Abstände müssen übereinstimmen.
\emph{Konsequenz:} $f$ darf Lichtkegel nicht verbiegen und muss Abstände $ds^2$ erhalten:
$ds'^2 = ds^2$. Das ist genau die Bedingung $f\in \mathrm{Iso}(\R^{1,3},\eta)$.

\textbf{A4. Empirische Forderung 3 (Gleiche Struktur der Gesetze).}
Freiheit (Trägheitsbewegung geradlinig–gleichförmig), Interferenzfiguren, Zerfallsraten, 
Maxwell- und Dirac-Gleichungen sollen \emph{im selben funktionalen Muster} erscheinen.
\emph{Konsequenz:} Wenn $S$ für eine Größe $\phi$ eine Gesetzesbeziehung „$D(\phi)=0$“ notiert,
muss $S'$ für die entsprechend übersetzte Größe $\phi'$ \emph{dieselbe Gleichungsform} haben.

\bigskip
\emph{B. Modellierungsprinzipien: Wie kodiert man die Übersetzung präzise?}

\textbf{B1. Raumzeit als Ereignismenge.}
Wir modellieren die Menge der Ereignisse als Minkowski-Raum $(M,g)=(\R^{1,3},\eta)$.
Ein Beobachterwechsel ist eine Isometrie $f:M\to M$ (Translation $+$ Lorentz-Abbildung).

\textbf{B2. Physikalische Größen als Funktionen/Felder.}
Skalare Größen (Temperatur, Wahrscheinlichkeit, Phase) sind Abbildungen $\phi:M\to\C$.
Vektor-/Spinorgrößen werden später durch zusätzliche lineare Strukturen modelliert.
\emph{Beobachter-$S$-Darstellung:} $\phi^{[S]}$.
\emph{Beobachter-$S'$-Darstellung:} $\phi^{[S']}$.

\textbf{B3. Übersetzungsoperator (Pullback) für Daten.}
„Dasselbe physikalische Objekt aus Sicht $S'$“ ist die \emph{umgezogene} Größe
\[
\phi' := f^*\phi := \phi\circ f.
\]
Dieser einfache Zug-Operator (Pullback) implementiert die Regel: 
„Bewerte die alte Größe am zugehörigen Ereignis“.

\textbf{B4. Gesetze als Operatorgleichungen.}
Ein Gesetz wird als Operatorgleichung $D(\phi)=0$ notiert. 
Beispiele:
$D=\Box_\eta+m^2$ (Klein–Gordon), $D=i\slashed{\partial}-m$ (Dirac), $D=\mathrm{d}^*\mathrm{d}$ (Maxwell in Potentialform).
$D$ ist \emph{lokal} (differentiell) und hängt von $g$ nur über $g$-kompatible Ableitungen ab.

\bigskip
\emph{C. Warum muss $f^*\!\circ D = D\circ f^*$ gelten? (Notwendigkeit aus den Axiomen)}

\textbf{C1. Operationales Kriterium „gleiche Vorhersagen“.}
Wenn $S$ eine Lösung $\phi$ (Messdaten/Initialdaten $\to$ Vorhersage) besitzt, muss $S'$ für dieselbe Versuchsanordnung die \emph{korrespondierende} Lösung $\phi':=f^*\phi$ besitzen. 
Andernfalls könnten $S$ und $S'$ in Widerspruch über beobachtbare Ereignisse geraten.

\textbf{C2. Forminvarianz als \underline{minimale} Strukturforderung.}
Die schwächste präzise Forderung, die das leistet, ist Kommutation:
\[
D(\phi)=0 \ \Longrightarrow\ D(f^*\phi)=0
\qquad\Longleftrightarrow\qquad
f^*\circ D = D\circ f^*.
\]
(Die Umkehrung ist ebenso notwendig, da $f$ invertierbar ist.)
Damit werden \emph{Gesetzesform} und \emph{Lösungsmenge} zwischen Beobachtern identifiziert.

\textbf{C3. Warum genau Isometrien?}
Nur wenn $f$ die Metrik erhält ($f^*g=g$), bleiben Lichtkegel und zeit-/raum-/lichtartige Klassen unverändert; andernfalls würden elementare Messprotokolle (Lichtuhren, Radarentfernung, Kausalordnung) zwischen Beobachtern nicht konsistent sein. 
Daraus folgt: $f\in\mathrm{Iso}(M,g)$ ist die maximale Verträglichkeitsklasse.

\bigskip
\emph{D. Konkrete Mikroszenarien, die die Forderung erzwingen}

\textbf{D1. Lichtuhr und Invarianz der Lichtgeschwindigkeit.}
Ein Lichtpuls zwischen zwei Spiegeln definiert eine Zeitnormale $\Delta\tau$. 
Zwei Beobachter müssen übereinstimmen, ob die Pulsbahn lichtartig ist. 
Dies erzwingt $f^*g=g$ und damit die Lorentzform von $f$.

\textbf{D2. Interferenzmuster eines Doppelspalts.}
Die Intensität $I(x)$ am Schirm ist eine Funktion der Phase $\Phi(x)$, die entlang $f$ umgezogen wird:
$I'(x')=I(f(x'))$. 
Das Muster (Nullstellen, Maxima) darf nicht vom Beobachter abhängen $\Rightarrow$ 
Gleichungsform (Wellengleichung/Phasenakkumulation) muss stabil sein $\Rightarrow$ 
$D$ kommutiert mit $f^*$.

\textbf{D3. Zerfallsstatistik eines instabilen Teilchens.}
Die Zerfallswahrscheinlichkeit $P(\text{„noch nicht zerfallen“};\tau)$ hängt von der Eigenzeit ab; 
deren Bestimmung nutzt $g$. 
In verschiedenen Inertialsystemen müssen dieselben Lebensdauergesetze gelten; 
nur die Darstellung (Zeitdilatation) ändert sich. 
Erneut folgt die Notwendigkeit der Kommutation.

\bigskip
\emph{E. Mathematische Verdichtung: Natürlichkeit von Operatoren}

\textbf{E1. Operatoren als geometrische Bauteile.}
Die Gesetzesoperatoren werden aus $g$ und der zugehörigen Levi-Civita-Verbindung gebaut 
(Laplace–Beltrami, D’Alembertian, Dirac via Spinverbindung). 
Diese Konstruktionen sind \emph{funktoriell}: 
zieht man Geometrie und Felder via $f^*$ um, ziehen sich auch die Operatoren mit.

\textbf{E2. Natürlichkeit.}
Formal heißt das
\[
f^*\big(D_g(\phi)\big) = D_{f^*g}\big(f^*\phi\big).
\]
Für Isometrien $f^*g=g$ reduziert sich dies zu
\[
\boxed{\quad f^*\circ D_g = D_g\circ f^* \quad}
\]
— genau unsere Forderung.

\bigskip
\emph{F. Zwei elementare Rechenchecks (Heuristik statt Tensorik)}

\textbf{F1. Klein–Gordon.}
$D=\Box_\eta+m^2$, $\phi'(x'):=\phi(\Lambda^{-1}x')$. 
Kettenregel $\partial'_\mu=(\Lambda^{-1})^\nu{}_\mu\partial_\nu$ und $\Lambda^T\eta\Lambda=\eta$ liefern
\[
(\Box'+m^2)\phi'(x') = \big((\Box+m^2)\phi\big)(\Lambda^{-1}x')=0.
\]

\textbf{F2. Dirac.}
$D=i\gamma^\mu\partial_\mu-m$, $\psi'(x')=S(\Lambda)\psi(\Lambda^{-1}x')$, 
$S^{-1}\gamma^\mu S=\Lambda^\mu{}_\nu\gamma^\nu$ liefert
\[
(i\gamma^\mu\partial'_\mu-m)\psi' = S(\Lambda)\,(i\gamma^\nu\partial_\nu-m)\psi\circ\Lambda^{-1}=0.
\]

\bigskip
\emph{G. Fazit (Logik der Konstruktion in drei Zeilen)}
\begin{enumerate}
\item Empirische Anforderungen (Vergleichbarkeit, Lichtkonstanz, gleiche Gesetzesform)
$\Longrightarrow$ Beobachterwechsel $f$ ist Isometrie und muss Lösungen auf Lösungen schicken.
\item Kodierung „gleicher physikalischer Zustand unter $f$“ als Pullback: $\phi\mapsto f^*\phi$.
\item Minimalpräzision der Gleichungsform: \emph{Kommutation mit Pullback}
\[
\boxed{\forall f\in\mathrm{Iso}(\R^{1,3},\eta):\quad f^*\circ D = D\circ f^*,\qquad
D(\phi)=0 \iff D(f^*\phi)=0.}
\]
\end{enumerate}

\bigskip

\emph{H. Was ist daran „Theorie“? – Der theoretische Charakter der SRT.}

Die Spezielle Relativitätstheorie ist keine einzelne Dynamik, sondern eine 
\emph{Meta-Theorie}, d.\,h. ein konsistentes geometrisches Rahmenwerk, das 
die Form aller physikalisch zulässigen Gleichungen einschränkt. 
Ihr „Theorie“-Charakter besteht also darin, dass sie definiert, 
\emph{welche Strukturen ein physikalisches Gesetz überhaupt haben darf}, 
damit es empirisch konsistent beschreibbar ist.  

\begin{enumerate}
  \item \textbf{Geometrischer Träger.}
  Die SRT postuliert, dass alle physikalischen Prozesse auf einer vierdimensionalen
  glatten Mannigfaltigkeit $M=\R^{1,3}$ stattfinden, ausgestattet mit der 
  Minkowski-Metrik $\eta$. 
  Diese Struktur ist nicht dynamisch: sie ist der feste Hintergrund, 
  auf dem Dynamiken ablaufen. 
  Das entspricht der empirischen Einsicht, dass Raum und Zeit nicht getrennt, 
  sondern untrennbar als einheitliche Raumzeit verstanden werden müssen.

  \item \textbf{Symmetriegruppe.}
  Die Menge der zulässigen Bezugssystemwechsel ist die Isometriegruppe
  $\mathrm{Iso}(M,\eta)=\mathcal P=\R^{1,3}\rtimes SO^+(1,3)$.
  Diese Gruppe ist physikalisch beobachtbar: sie beschreibt die Menge aller möglichen
  Inertialsysteme, zwischen denen experimentell kein Unterschied gemacht werden kann.
  Ihre Wirkung auf physikalische Felder modelliert die Transformation von Beschreibungen
  zwischen Beobachtern.

  \item \textbf{Funktionsprinzip.}
  Die Übersetzungsvorschrift $f^*$ zwischen Beobachterdarstellungen
  ist keine Zusatzannahme, sondern eine formale Konsequenz aus der Forderung,
  dass physikalische Größen geometrische Objekte auf $(M,\eta)$ sind.
  Der Pullback $f^*$ ist der natürliche Funktor, der „dasselbe physikalische Objekt“
  in unterschiedlichen Koordinatendarstellungen ausdrückt.

  \item \textbf{Kriterium theoretischer Zulässigkeit.}
  Eine Gleichung $D(\phi)=0$ beschreibt genau dann ein physikalisches Gesetz im Sinn der SRT,
  wenn sie \emph{natürlich} unter allen Isometrien ist, also
  \[
  f^*\circ D = D\circ f^*.
  \]
  Dies ist die präzise mathematische Übersetzung des Relativitätsprinzips:
  die Form der Gleichung ist unabhängig vom gewählten Inertialsystem.

  \item \textbf{Konkrete Theorien als Realisierungen.}
  Einzelne physikalische Theorien (Klein–Gordon, Maxwell, Dirac, Yang–Mills)
  sind dann Spezialisierungen dieses Rahmens:
  \[
  D_{\text{KG}}=\Box_\eta+m^2,\qquad 
  D_{\text{Max}}=\mathrm{d}^*\mathrm{d},\qquad
  D_{\text{Dir}}=i\gamma^\mu\partial_\mu-m,\quad\text{etc.}
  \]
  Jede erfüllt die Bedingung $f^*\circ D = D\circ f^*$ für alle $f\in\mathcal P$.
  Ihre Verschiedenheit liegt nur in der Wahl des Feldbündels $\mathcal F$
  (Skalar-, Vektor-, Spinorbündel) und der daraus konstruierten natürlichen Operatoren.

  \item \textbf{Theorie als geometrisches Postulat.}
  In diesem Sinn ist die SRT eine \emph{geometrische Theoriekategorie}:
  sie postuliert, dass „Physik = Geometrie“ auf $(M,\eta)$ ist.
  Jede zulässige physikalische Größe ist ein geometrischer Abschnitt eines
  Bündels über $(M,\eta)$, jede zulässige Gleichung ein natürlicher Differentialoperator.
  Damit reduziert sich die empirische Forderung „gleiche Physik in allen Bezugssystemen“
  auf eine rein formale Invarianzforderung.
\end{enumerate}

\textbf{Zusammenfassung.}
\begin{itemize}
  \item Die SRT liefert den \emph{geometrischen Rahmen}, nicht die Dynamik.
  \item Sie definiert, was ein „Bezugssystem“ und ein „Gesetz“ überhaupt formal bedeutet.
  \item Sie garantiert die Vergleichbarkeit physikalischer Aussagen zwischen Beobachtern.
  \item Dynamische Theorien (Maxwell, Dirac, etc.) sind \emph{Modelle innerhalb dieses Rahmens}.
\end{itemize}

Damit ist die Spezielle Relativitätstheorie die mathematische Formalisierung
des Relativitätsprinzips selbst — eine Theorie \emph{zweiter Ordnung}, 
deren Inhalt nicht ein bestimmtes Gesetz, sondern der Begriff von Gesetzmäßigkeit ist.


\pagebreak


\subsection{Bezugssystemwechsel als Rahmenwechsel im Minkowski-Raum}

\emph{Motivation.}
Ein physikalisches \emph{Bezugssystem} (Inertialsystem) ist konzeptionell eine Wahl von 
Zeit- und Raumrichtungen, mit denen man physikalische Größen als Komponenten beschreibt.
Differentialgeometrisch entspricht dies einer Wahl eines \emph{orthonormalen Rahmens} 
auf der Raumzeitmannigfaltigkeit.  
Ein Wechsel des Bezugssystems wird dann zu einem \emph{Rahmenwechsel}, der durch eine
Lorentztransformation beschrieben wird.  
In diesem Sinn ist der Begriff „Beobachterwechsel“ äquivalent zur 
Wahl eines anderen lokalen Rahmens im Tangentialraum.

\bigskip
\emph{1. Orthonormale Rahmen.}

Sei $(M,g)$ eine vierdimensionale Lorentz-Mannigfaltigkeit.
Ein \emph{(lokaler) Rahmen} am Punkt $p\in M$ ist eine geordnete Basis
\[
e(p)=(e_0(p),e_1(p),e_2(p),e_3(p))
\]
des Tangentialraums $T_pM$.
Der Rahmen heißt \emph{orthonormal}, wenn gilt
\[
g(e_\mu,e_\nu) = \eta_{\mu\nu} = \mathrm{diag}(1,-1,-1,-1).
\]
Solche Rahmen bilden das \emph{Lorentz-Rahmenbündel}
\[
\mathrm{Fr}_\eta(M)
= \{ (p,e) \mid e=(e_0,\ldots,e_3) \text{ orthonormal bzgl. } g \}.
\]
Dieses ist ein \emph{Hauptbündel} mit Strukturgruppe $O(1,3)$.

\bigskip
\emph{2. Gruppenwirkung und Beobachterwechsel.}

Die Lorentzgruppe $O(1,3)$ wirkt \emph{rechts} auf $\mathrm{Fr}_\eta(M)$ durch
\[
e \cdot \Lambda := (e'_\mu), \qquad e'_\mu := e_\nu \Lambda^\nu{}_\mu.
\]
Dies ist eine freie und transitive Wirkung auf jeder Faser:
alle orthonormalen Rahmen in $T_pM$ sind durch eine eindeutige
Lorentztransformation ineinander überführbar.

\textbf{Interpretation:}
\begin{itemize}
  \item Ein bestimmter Rahmen $e=(e_0,e_1,e_2,e_3)$ entspricht einem 
  Inertialsystem (Beobachter) mit Zeitrichtung $e_0$ und Raumbasis $e_i$.
  \item Ein neues Bezugssystem wird durch Anwendung einer Lorentztransformation 
  $\Lambda$ beschrieben: $e\mapsto e\cdot\Lambda$.
  \item In der Sprache der Physik: dieselben Ereignisse werden 
  in einem \emph{verkippten} Zeit-/Raumraster dargestellt.
\end{itemize}

\bigskip
\emph{3. Der Fall des Minkowski-Raums.}

Sei $(M,g)=(\R^{1,3},\eta)$ mit $\eta=\mathrm{diag}(1,-1,-1,-1)$.
Dann ist jeder Rahmen global trivial, und das Rahmenbündel ist
\[
\mathrm{Fr}_\eta(\R^{1,3}) \cong \R^{1,3}\times O(1,3).
\]

\textbf{Standardrahmen:}
\[
e = (\partial_t, \partial_x, \partial_y, \partial_z),
\qquad
g(e_\mu,e_\nu)=\eta_{\mu\nu}.
\]

\textbf{Lorentz-Boost in $z$-Richtung:}
Für eine Rapidität $\eta$ (mit $\tanh\eta=v$) sei
\[
\Lambda(\eta)
=
\begin{pmatrix}
\cosh\eta & 0 & 0 & -\sinh\eta\\
0 & 1 & 0 & 0\\
0 & 0 & 1 & 0\\
-\sinh\eta & 0 & 0 & \cosh\eta
\end{pmatrix}.
\]
Dann ist der neue Rahmen
\[
e' = e \cdot \Lambda(\eta),
\]
d.\,h.
\[
\begin{aligned}
e'_0 &= \cosh\eta\,\partial_t - \sinh\eta\,\partial_z,\\
e'_3 &= -\sinh\eta\,\partial_t + \cosh\eta\,\partial_z,\\
e'_1 &= \partial_x,\quad e'_2=\partial_y.
\end{aligned}
\]
Man überprüft direkt:
\[
g(e'_\mu,e'_\nu)=\eta_{\mu\nu}.
\]
Physikalisch beschreibt $e'_0$ den Vierergeschwindigkeitsvektor eines Beobachters,
der sich mit Geschwindigkeit $v=\tanh\eta$ entlang der $z$-Achse bewegt.

\bigskip
\emph{4. Komponentenwechsel physikalischer Größen.}

Sei $T$ ein Tensorfeld auf $M$, dessen Komponenten im Rahmen $e$ gegeben sind durch
\[
T = T^{\mu_1\ldots}_{\nu_1\ldots}\, e_{\mu_1}\otimes\cdots\otimes e^{\nu_k}.
\]
Unter dem Rahmenwechsel $e\mapsto e'=e\cdot\Lambda$ transformieren die Komponenten
nach der bekannten Regel:
\[
T'^{\mu_1\ldots}_{\nu_1\ldots}
= 
\Lambda^{\mu_1}{}_{\rho_1}\cdots
\Lambda_{\nu_1}{}^{\sigma_1}\cdots
T^{\rho_1\ldots}_{\sigma_1\ldots}.
\]
Somit wird der „Bezugssystemwechsel“ im physikalischen Sinn formal als 
\emph{lineare Gruppenwirkung der Lorentzgruppe auf die Komponenten im Rahmenbündel}
realisiert.

\bigskip
\emph{5. Physikalische Interpretation.}

\begin{itemize}
  \item Ein \emph{Beobachter} wird durch einen zeitartigen Einheitsvektor $u(p)$
  beschrieben (der $e_0$-Richtung des Rahmens entspricht).
  \item Ein \emph{Bezugssystem} ist die Erweiterung dieses Vektors zu einem
  vollständigen orthonormalen Rahmen $(e_0,\ldots,e_3)$.
  \item Ein \emph{Wechsel des Bezugssystems} ist eine Transformation 
  $e\mapsto e\cdot\Lambda$, bei der sich sowohl Zeit- als auch Raumachsen
  hyperbolisch „verkippen“ (Lorentz-Boost).
  \item Physikalische Gleichungen sind \emph{invariant} unter dieser Gruppenwirkung,
  d.\,h. die Operatoren $D$ erfüllen $f^*\circ D = D\circ f^*$, 
  wobei $f$ die durch $\Lambda$ induzierte Isometrie ist.
\end{itemize}

\bigskip
\emph{6. Fazit.}
\begin{itemize}
  \item Der \textbf{Bezugssystemwechsel} ist geometrisch identisch mit einem
  \textbf{Rahmenwechsel} im orthonormalen Rahmenbündel $\mathrm{Fr}_\eta(M)$.
  \item Die \textbf{Lorentztransformationen} sind genau die 
  \textbf{Strukturgruppenoperationen} dieses Bündels.
  \item Die \textbf{Forminvarianz physikalischer Gesetze} 
  entspricht der \textbf{Natürlichkeit} der zugehörigen Operatoren unter dieser Gruppenwirkung.
\end{itemize}

Somit erhält der intuitive Begriff „Beobachterwechsel“ seine mathematisch präzise 
Gestalt als \emph{Hauptbündelaktion} der Lorentzgruppe auf den orthonormalen Rahmen 
des Minkowski-Raums.

\subsection{Bündeltheorie $\to$ „Lorentztrafos in der Praxis“ (Top–Down)}

\emph{Ziel.} Wir geben eine bündeltheoretische Formulierung der SRT und leiten daraus
in einem \emph{Top–Down}-Sinn die üblichen Transformationsregeln im Minkowski-Raum
$\R^{1,3}$ her. Leitidee: \emph{Bezugssystemwechsel} = \emph{Rahmenwechsel} im
orthonormalen Rahmenbündel; \emph{Felder} = \emph{Schnitte assoziierter Bündel};
\emph{Lorentztransformation} = \emph{rechte Hauptbündelaktion}.

\bigskip
\emph{1. Hauptbündel der orthonormalen Rahmen.}

Sei $(M,g)$ eine Lorentz-Mannigfaltigkeit (in SRT: $M=\R^{1,3}$, $g=\eta$).
Das \emph{Lorentz-Rahmenbündel}
\[
\mathrm{Fr}_\eta(M)\;:=\;\bigl\{(p,e)\mid p\in M,\ e=(e_0,\dots,e_3)\text{ Basis von }T_pM,\ g(e_\mu,e_\nu)=\eta_{\mu\nu}\bigr\}
\]
ist ein \emph{Hauptbündel} mit Strukturgruppe $O(1,3)$:
\[
\mathrm{Fr}_\eta(M)\xrightarrow{\ \pi\ }M,\qquad
R_\Lambda:\mathrm{Fr}_\eta(M)\to\mathrm{Fr}_\eta(M),\quad (p,e)\mapsto (p,e\cdot\Lambda),\quad \Lambda\in O(1,3).
\]
Die zusammenhängende, orientierungs- und zeitorientierungserhaltende Komponente ist $SO^+(1,3)$.

\smallskip
\emph{Physik:} Ein Punkt $e\in\mathrm{Fr}_\eta(M)$ ist ein \emph{Beobachterrahmen} am Ereignis $p$.
Die rechte Gruppenwirkung $R_\Lambda$ realisiert den \emph{Bezugssystemwechsel}.

\bigskip
\emph{2. Assoziierte Bündel und Felder als Schnitte.}

Sei $\rho:G\to\GL(V)$ eine (endlichdimensionale) Darstellung von $G:=SO^+(1,3)$
(z.\,B. die triviale Darstellung für Skalare, die Standarddarstellung für Vektoren,
oder die Spinor-Darstellung nach Hebung auf $\Spin(1,3)$).
Dann ist das \emph{assoziierte Vektorbündel}
\[
E_\rho \;:=\; \mathrm{Fr}_\eta(M)\times_\rho V\ \ \longrightarrow\ \ M
\]
über die Äquivalenz $(e,v)\sim (e\Lambda,\rho(\Lambda)^{-1}v)$ definiert.
Ein \emph{Feld} vom Typ $\rho$ ist ein glatter Schnitt $\phi\in\Gamma(E_\rho)$.

\smallskip
\emph{Physik:} „Ein skalares Feld“ $\phi$ ist $\rho=\mathbf{1}$; „Vierervektorfeld“ ist
$\rho=\text{Std}$; \emph{Spinorfeld} benötigt eine Spin-Struktur
$\widetilde{\mathrm{Fr}}_\eta(M)\to\mathrm{Fr}_\eta(M)$ und eine Darstellung $\rho:\Spin(1,3)\to\GL(\mathbb{C}^4)$.

\bigskip
\emph{3. Lokale Darstellungen und Eichsicht („lokale Lorentzsymmetrie“).}

Eine lokale \emph{Rahmensektion} $s:U\subset M\to\mathrm{Fr}_\eta(M)$ (ein \emph{Vierbein/Tetrad})
wählt in jedem Punkt $p\in U$ einen orthonormalen Rahmen $e_\mu(p)$.
Über $s$ identifiziert man $E_\rho|_U \cong U\times V$ und erhält Komponenten
\[
\phi^{[s]}:U\to V,\qquad \text{„Feld in der Darstellung relativ zum Rahmen }s\text{“}.
\]
Wechselt man den Rahmen lokal durch $s' = s\cdot\Lambda$ (mit glattem $\Lambda:U\to SO^+(1,3)$),
so transformiert die Komponentendarstellung \emph{punktweise} durch die Faseraktion:
\[
\phi^{[s']}(p)\;=\;\rho\bigl(\Lambda(p)\bigr)\,\phi^{[s]}(p).
\]
Dies ist genau die \emph{(lokale) Lorentz-„Eichtransformation“}.

\smallskip
\emph{Physik:} Komponenten eines Feldes hängen von der Rahmenwahl ab; die \emph{Geometrie} des Felds
(i.\,e. der Schnitt) ist rahmenunabhängig. „Beobachterwechsel“ ist die \emph{Gleichheit des Schnitts}
bei veränderter Komponentendarstellung.

\bigskip
\emph{4. Verbindungen und kovariante Ableitungen.}

Eine \emph{(Ehresmann-)Verbindung} auf $\mathrm{Fr}_\eta(M)$ ist eine $ \mathfrak{so}(1,3)$-wertige 1-Form
$\omega$ auf $\mathrm{Fr}_\eta(M)$ mit den üblichen Äquivarianzeigenschaften. Für einen lokalen Rahmen $s$
ist die \emph{lokale Verbindungsform} $\omega^{[s]}:=s^*\omega\in\Omega^1\bigl(U,\mathfrak{so}(1,3)\bigr)$.
Die zu $\rho$ induzierte \emph{kovariante Ableitung} auf $E_\rho$ hat in Komponenten die Gestalt
\[
\nabla^{[s]} \phi^{[s]}\;=\; d\phi^{[s]} \;+\; \rho_*\bigl(\omega^{[s]}\bigr)\,\phi^{[s]},
\quad
\rho_*:\mathfrak{so}(1,3)\to\mathfrak{gl}(V).
\]
Unter einem lokalen Rahmenwechsel $s\mapsto s\cdot\Lambda$ transformieren
\[
\begin{aligned}
\phi^{[s]} &\mapsto \phi^{[s']}\;=\;\rho(\Lambda)\,\phi^{[s]},\\
\omega^{[s]} &\mapsto \omega^{[s']} \;=\; \Lambda^{-1}\omega^{[s]}\Lambda + \Lambda^{-1}d\Lambda,
\end{aligned}
\]
und \emph{kompatibel} dazu
\[
\nabla^{[s']} \phi^{[s']} \;=\; \rho(\Lambda)\,\bigl(\nabla^{[s]} \phi^{[s]}\bigr).
\]
Damit ist die kovariante Ableitung ein \emph{natürlicher Operator}.

\smallskip
\emph{Physik:} In SR (Minkowski, global inertial) kann man $s$ global so wählen, dass
$\omega^{[s]}=0$ (flache, torsionsfreie Levi-Civita-Verbindung, globale Lorentz-Rahmen):
dann reduziert sich $\nabla$ auf \emph{partielle Ableitungen} (für Tensorfelder).
Für Spinorfelder verbleibt bei Wahl eines global konstanten Vierbeins ebenfalls $\omega=0$,
sodass die \emph{Spinverbindung} in kartesischen Inertialkoordinaten verschwindet
(Dirac-Operator $i\gamma^\mu\partial_\mu$).

\bigskip
\emph{5. Natürliche Differentialoperatoren und Lorentz-Kovarianz.}

Ein (lokal definierter) Differentialoperator
\[
D:\Gamma(E_\rho)\longrightarrow \Gamma(E_\sigma)
\]
heißt \emph{natürlich (Lorentz-kovariant)}, wenn für jede Isometrie
$f\in\Iso(M,g)$ und die induzierte Bündelabbildung $f^*:E_\rho\to E_\rho$ gilt:
\[
f^*\circ D \;=\; D\circ f^*.
\]
Lokal in einem Rahmen $s$ liest sich das als \emph{Rahmenwechsel-Kovarianz}:
wechselt man $s\mapsto s\cdot\Lambda$, so gehen die Komponenten von $D$ durch die
darstellungsinduzierte Konjugation über (die Operatorform bleibt gleich).

\smallskip
\emph{Physik:} Genau dies ist die präzise Fassung der „Forminvarianz der Gesetze“.
Beispiele (SR, flach, globaler Inertialrahmen):
\[
\begin{aligned}
&\text{Klein–Gordon:}& D_{\text{KG}} &= \Box_\eta + m^2,\\[2pt]
&\text{Maxwell (in Potentialform):}& D_{\text{Max}} &= d^\dagger d \quad \text{auf 1-Formen},\\[2pt]
&\text{Dirac:}& D_{\text{Dir}} &= i\gamma^\mu\partial_\mu - m.
\end{aligned}
\]
Unter einem globalen Lorentzboost $\Lambda\in SO^+(1,3)$ wirken die Felder über $\rho(\Lambda)$
und die Ableitungen gemäß Kettenregel; wegen $\Lambda^T\eta\Lambda=\eta$ bleibt die Form der
Operatoren erhalten: $f^*D = D f^*$.

\bigskip
\emph{6. Spezialisierung: Minkowski-Raum als trivialer Fall.}
Für $(M,g)=(\R^{1,3},\eta)$ ist $\mathrm{Fr}_\eta(M)\cong \R^{1,3}\times SO^+(1,3)$ \emph{trivial}.
Wählt man eine globale Sektion $s$ (z.\,B. den Standardrahmen $e_\mu=\partial_\mu$),
so sind:
\begin{itemize}
\item alle assoziierten Bündel \emph{global trivial}, $E_\rho\cong \R^{1,3}\times V$,
\item die Levi-Civita-Verbindung ist \emph{flach und inerten}: $\omega^{[s]}=0$,
\item kovariante Ableitungen sind \emph{partielle Ableitungen} (für Tensor- und Spinorfelder),
\item globale Lorentztrafos $\Lambda$ handeln \emph{konstant} in der Faser: 
$\phi(x)\mapsto \rho(\Lambda)\,\phi(\Lambda^{-1}x)$.
\end{itemize}
Damit reduziert sich der bündeltheoretische Apparat exakt auf die vertrauten
Regeln der SRT: „Feld transformiert mit einer (Spin-)Darstellung, Argument mit
der inversen Lorentzmatrix“.

\bigskip
\emph{7. „Physikalischer Stil“ – Arbeitsrezept aus der Bündelsicht.}
\begin{enumerate}
\item \textbf{Feldtyp wählen:} Wähle eine Darstellungsart $\rho$ von $SO^+(1,3)$ 
(oder von $\Spin(1,3)$ für Spinoren).
\item \textbf{Bündel und Rahmen:} Identifiziere $E_\rho=\mathrm{Fr}_\eta(M)\times_\rho V$. 
Fixiere einen globalen Inertialrahmen $s$ (Minkowski).
\item \textbf{Operatorbau:} Baue $D$ aus $g$ und $\nabla$ (hier: $\partial$) \emph{representationstreu} 
(d.\,h. alle Indizes mit $\eta$ heben/senken, $\gamma^\mu$ für Spinoren etc.).
\item \textbf{Lorentztest (Kovarianz):} Prüfe $f^*D = D f^*$ für $f$ aus $\Iso(M,\eta)$
(äquivalent: Forminvarianz der Komponenten unter $s\mapsto s\cdot\Lambda$).
\item \textbf{Rechnen:} Nutze globale Trivialität: Komponentenrechnung in $s$,
Lorentzboost via konstante $\Lambda$ und $\rho(\Lambda)$; Kettenregel für Argumente.
\end{enumerate}

\bigskip
\emph{8. Zwei explizite Mikrotests.}

\emph{(i) Klein–Gordon.} $\phi:\R^{1,3}\to\C$, $\rho=\mathbf{1}$.
Unter $x'=\Lambda x$ und $\phi'(x')=\phi(\Lambda^{-1}x')$:
\[
(\Box'+m^2)\phi'(x') \;=\; (\Box+m^2)\phi(\Lambda^{-1}x') \;=\; 0,
\quad \text{da }\Lambda^T\eta\Lambda=\eta.
\]
\emph{(ii) Dirac.} $\psi:\R^{1,3}\to\C^4$, $\rho=S:\Spin(1,3)\to\GL(\C^4)$ mit
$S^{-1}\gamma^\mu S=\Lambda^\mu{}_\nu\gamma^\nu$.
Setze $\psi'(x')=S(\Lambda)\psi(\Lambda^{-1}x')$:
\[
(i\gamma^\mu\partial'_\mu-m)\psi' \;=\; S(\Lambda)\,(i\gamma^\nu\partial_\nu-m)\psi\;\circ\Lambda^{-1}\;=\;0.
\]
Beide Tests sind Sichtbarmachungen von $f^*D=Df^*$ im trivialen Bündelfall.

\bigskip
\emph{9. Warum diese Sicht „mehr ist als Koordinatenrechnung“.}

Die Bündel- und Hauptbündelsicht
\begin{itemize}
\item erklärt \emph{konzeptuell}, warum Komponentenformeln so (und nur so) transformieren,
\item trennt \emph{Geometrie} (Schnitt, natürliche Operatoren) von \emph{Darstellung}
(Rahmen-/Eichwahl),
\item generalisiert reibungslos zu gekrümmten Hintergründen (GR): 
$e^a{}_\mu$ (Vierbein), $\omega^a{}_b$ (Spinverbindung), 
lokale Lorentzsymmetrie als \emph{innere Eichgruppe} des Rahmenbündels,
\item macht den \emph{Theorieinhalt} „Forminvarianz/Natürlichkeit“ präzise
und unabhängig von Koordinaten.
\end{itemize}

\bigskip
\emph{10. Quintessenz.}

\emph{Top–Down} vom Hauptbündel $\mathrm{Fr}_\eta(M)$ über assoziierte Bündel $E_\rho$
zu natürlichen Operatoren $D$ ist der präzise mathematische Kern hinter der
physikalischen Aussage „Lorentzkovarianz“. Im Minkowski-Raum führt die
Trivialität des Bündels direkt zu den gewohnten Transformationsregeln:
\[
\boxed{\ \phi'(x')=\rho(\Lambda)\,\phi(\Lambda^{-1}x'),\qquad
D' = D,\ \ f^*D=Df^*.\ }
\]
Diese Formulierung ist nicht nur „schön“, sondern liefert die exakte Brücke
zwischen \emph{Bezugssystemwechsel} (Rahmenwechsel) und \emph{praktischer
Rechnung} im physikalischen Alltag.


%–– Klein–Gordon im Minkowski-Rahmen: Felder als Schnitte, Operator, Lösungen, Transformationen ––%

\subsection*{Klein–Gordon im SRT-Rahmen – formale Fassung des physikalischen Jargons}

\emph{Raumzeit.}
Sei die Minkowski-Raumzeit
\[
(M,g)=(\R^{1,3},\eta),\qquad 
\eta=\mathrm{diag}(1,-1,-1,-1),
\]
mit globalen Inertialkoordinaten $x=(x^\mu)=(t,\mathbf x)$ und $\partial_\mu=\frac{\partial}{\partial x^\mu}$.
Die Isometriegruppe ist die Poincaré-Gruppe
\[
\mathcal P=\R^{1,3}\rtimes SO^+(1,3)=:\{(a,\Lambda)\},
\]
mit Wirkung auf $M$ gegeben durch $(a,\Lambda)\cdot x := \Lambda x + a$.

\emph{Feld als Schnitt.}
Ein (komplexes) skalares Klein–Gordon-Feld ist ein glatter Schnitt des trivialen Linienbündels
\[
E:=M\times\C \;\longrightarrow\; M,\qquad 
\Gamma(E)\cong C^\infty(M,\C).
\]
Physikalischer Jargon „\emph{$\phi$ ist ein (Skalar-)Feld}“ entspricht formal
\[
\phi\in \Gamma(E).
\]

\emph{D’Alembert-Operator und Klein–Gordon-Operator.}
Mit $\eta^{\mu\nu}$ der inversen Metrik ist der D’Alembert-Operator
\[
\square_\eta \;:=\; \eta^{\mu\nu}\,\partial_\mu\partial_\nu \;=\; \partial_t^2-\Delta_{\mathbf x}.
\]
Für eine feste Masse $m\ge 0$ ist der Klein–Gordon-Operator
\[
P_m \;:=\; \square_\eta + m^2 \;:\; C^\infty(M)\longrightarrow C^\infty(M).
\]
Jargon „\emph{der Klein–Gordon-Operator wirkt auf $\phi$}“ heißt formal: $P_m(\phi)$ wird punktweise
durch obige Differentialausdrücke berechnet.

\emph{Lösungen.}
„\emph{$\phi$ ist Lösung der Klein–Gordon-Gleichung}“ bedeutet formal:
\[
P_m(\phi)=0\qquad\Longleftrightarrow\qquad \phi\in\ker P_m \;=:\; \mathrm{Sol}_m^{\infty}.
\]
Je nach Regularitätsklasse nutzt man dieselbe Definition in verallgemeinerten Räumen:
\[
\mathrm{Sol}_m^{\mathrm{test}}:=\{\phi\in C_c^\infty(M):P_m\phi=0\},\quad
\mathrm{Sol}_m^{\mathcal S}:=\{\phi\in\mathcal S(M):P_m\phi=0\},\quad
\mathrm{Sol}_m^{\mathcal D'}:=\{u\in\mathcal D'(M):P_m u=0\}.
\]
Im Cauchy-Bild sind klassische Lösungen durch Anfangsdaten auf $\Sigma=\{t=0\}\cong\R^3$ parametrisiert:
\[
\phi|_\Sigma=f,\qquad \partial_t\phi|_\Sigma=g,
\]
wobei $(f,g)$ geeignete Regularität (z.\,B. Schwarz- oder Sobolev-Daten) erfüllen; dann existiert eindeutig
$\phi\in \mathrm{Sol}_m^{\infty}$ mit diesen Daten.

\emph{Poincaré-Wirkung auf Schnitte (Bezugssystemwechsel).}
Die Isometrie $(a,\Lambda)\in\mathcal P$ induziert einen \emph{Pullback} auf Schnitten:
\[
\bigl((a,\Lambda)^*\phi\bigr)(x)
\;:=\; \phi\bigl((a,\Lambda)\cdot x\bigr)
\;=\; \phi(\Lambda x + a).
\]
Jargon „\emph{Rahmenwechsel/Boost/Translation wirkt auf das Feld}“ heißt formal:
\[
\phi \;\longmapsto\; (a,\Lambda)^*\phi.
\]
Für Ableitungen gilt die Kettenregel
\[
\partial_\mu\bigl((a,\Lambda)^*\phi\bigr)(x)
= \Lambda^\nu{}_{\mu}\, (\partial_\nu\phi)\bigl(\Lambda x+a\bigr),
\]
und daher
\[
\square_\eta\bigl((a,\Lambda)^*\phi\bigr)(x)
= \eta^{\mu\nu}\Lambda^\rho{}_{\mu}\Lambda^\sigma{}_{\nu}\,
(\partial_\rho\partial_\sigma\phi)\bigl(\Lambda x+a\bigr)
= \square_\eta \phi\bigl(\Lambda x+a\bigr),
\]
weil $\Lambda^T\eta\Lambda=\eta$. Also \emph{kommutiert} der Operator mit dem Pullback:
\[
\boxed{\quad (a,\Lambda)^*\circ P_m \;=\; P_m\circ (a,\Lambda)^* \quad}
\]
(„\emph{Lorentz/Poincaré-Kovarianz}“). Insbesondere ist der Lösungsraum invariant:
\[
\phi\in\mathrm{Sol}_m^\infty \;\Longrightarrow\; (a,\Lambda)^*\phi\in\mathrm{Sol}_m^\infty.
\]

\emph{Fourier-/Massenschalen-Charakterisierung (optional präzisiert).}
Für temperierte Lösungen $\phi\in \mathcal S'(M)$ gilt äquivalent
\[
P_m\phi=0 
\quad\Longleftrightarrow\quad
\widehat{\phi}(p)\ \text{ist auf der Massenschale } \{p\in\R^{1,3}:p^2=m^2\} \text{ getragen,}
\]
wobei $p^2=\eta^{\mu\nu}p_\mu p_\nu$ und die Fourier-Transformation mit Konvention
$\widehat{\phi}(p)=\int e^{ip\cdot x}\phi(x)\,dx$ genommen sei. 
Die Poincaré-Wirkung entspricht in Impulsraum der Darstellung
\[
\widehat{(a,\Lambda)^*\phi}(p) \;=\; e^{+i\,a\cdot p}\,\widehat{\phi}(\Lambda^{-1}p),
\]
die die Massenschale invariant lässt.

\emph{Physik $\leftrightarrow$ Formal – Mini-Wörterbuch.}
\begin{itemize}
  \item „\emph{Skalares Feld in der SRT}“ $\;\rightsquigarrow\;$ Schnitt $\phi\in\Gamma(E)\cong C^\infty(M)$.
  \item „\emph{Klein–Gordon-Gleichung}“ $\;\rightsquigarrow\;$ $P_m\phi=0$ mit $P_m=\square_\eta+m^2$.
  \item „\emph{Wechsel des Bezugssystems/Boost/Translation}“ $\;\rightsquigarrow\;$ Pullback $(a,\Lambda)^*$.
  \item „\emph{Gleichungsform bleibt gleich}“ $\;\rightsquigarrow\;$ $(a,\Lambda)^* P_m = P_m (a,\Lambda)^*$.
  \item „\emph{Lösung transformiert wieder zu Lösung}“ $\;\rightsquigarrow\;$ $(a,\Lambda)^*(\ker P_m)=\ker P_m$.
\end{itemize}

\emph{Zusammenfassung.}
Im SRT-Rahmen ist ein Klein–Gordon-Feld ein Schnitt des trivialen Linienbündels über $(\R^{1,3},\eta)$.
Der Klein–Gordon-Operator $P_m=\square_\eta+m^2$ ist ein natürlicher, von der Metrik induzierter
Differentialoperator. Lösungen sind genau die Schnitte im Kern von $P_m$; sie bilden einen
Poincaré-invarianten Unterraum von $\Gamma(E)$, da $P_m$ mit der Pullback-Wirkung der
Poincaré-Gruppe kommutiert.


\pagebreak

%–– Bezugssystemwechsel via Rahmenbündel & Schnitträume; Verhalten von Differentialoperatoren ––%

\subsection*{Bezugssystemwechsel aus der Bündel-/Schnitt-Perspektive}

\emph{Grunddaten.}
Sei $(M,g)$ eine Lorentz-Mannigfaltigkeit (in der SRT: $M=\R^{1,3}$, $g=\eta$).
Das \emph{Lorentz-Rahmenbündel} (Hauptbündel)
\[
\mathrm{Fr}_\eta(M):=\{(p,e)\mid p\in M,\ e=(e_0,\dots,e_3)\text{ Basis von }T_pM,\ g(e_\mu,e_\nu)=\eta_{\mu\nu}\}
\]
trägt die rechte $G$-Wirkung mit $G:=SO^+(1,3)$:
\[
R_\Lambda:\ (p,e)\mapsto (p,e\cdot\Lambda),\qquad (e\cdot\Lambda)_\mu:=e_\nu\,\Lambda^\nu{}_\mu.
\]
Für eine (endlichdimensionale) Darstellung $\rho:G\to\GL(V)$ sei das \emph{assoziierte Bündel}
\[
E_\rho:=\mathrm{Fr}_\eta(M)\times_\rho V\ \longrightarrow\ M,
\qquad (e,v)\sim (e\cdot\Lambda,\rho(\Lambda)^{-1}v).
\]
Ein \emph{Feld vom Typ $\rho$} ist ein glatter Schnitt $\phi\in\Gamma(E_\rho)$.

\emph{Schnitte als äquivariante Abbildungen.}
Jeder Schnitt $\phi\in\Gamma(E_\rho)$ entspricht eindeutig einer glatten Abbildung
\[
\Phi:\mathrm{Fr}_\eta(M)\to V,\qquad \Phi(e\cdot\Lambda)=\rho(\Lambda)^{-1}\,\Phi(e)
\quad\text{(Rechts-$G$-Äquivarianz),}
\]
via $\phi(p)=[e,\Phi(e)]$ (unabhängig von der Wahl $e\in\mathrm{Fr}_\eta(M)_p$).
Dies identifiziert den \emph{Raum der Schnitte} mit einem Äquivariantenraum:
\[
\Gamma(E_\rho)\ \cong\ C^\infty(\mathrm{Fr}_\eta(M),V)^{G\text{-eq}}.
\]

\emph{Bezugssystem (lokaler Rahmen) und seine Änderung.}
Eine (lokale) \emph{Rahmensektion} $s:U\to\mathrm{Fr}_\eta(M)$ (auch Vierbein/Tetrad) liefert
eine triviale Identifikation $E_\rho|_U\simeq U\times V$ über
\[
\phi\ \longleftrightarrow\ \phi^{[s]}:U\to V,\qquad
\phi(p)=[s(p),\,\phi^{[s]}(p)].
\]
Wechselt man den Rahmen lokal durch $s' = s\cdot\Lambda$ (mit glatter $\Lambda:U\to G$), so gilt
\[
\phi^{[s']}(p)\;=\;\rho\!\big(\Lambda(p)\big)\,\phi^{[s]}(p).
\]
Dies ist der \emph{passive} Komponentenwechsel („Bezugssystemwechsel“ als \emph{Rahmenwechsel}).

\emph{Isometrien als aktive Beobachterwechsel.}
Sei $f\in\Iso(M,g)$ eine Isometrie. Es existiert eine eindeutige liftende Bündelabbildung
\[
\tilde f:\mathrm{Fr}_\eta(M)\to \mathrm{Fr}_\eta(M),\qquad \tilde f(p,e)=(f(p),\ d f_p\circ e),
\]
mit $\tilde f\circ R_\Lambda=R_\Lambda\circ \tilde f$.
Auf Schnitten (als äquivariante $\Phi$) wirkt $f$ durch \emph{Pullback}
\[
(f^*_\rho\phi)\ \widehat{=}\ \Phi\circ \tilde f:\quad
\Phi_f(e):=\Phi\!\big(\tilde f(e)\big).
\]
Im Trivialisat $s$ liest sich das als
\[
(f^*_\rho\phi)^{[s]}(p)
=\rho\!\big(\Lambda_f(p)\big)\,\phi^{[s_f]}\!\big(f(p)\big),
\]
wobei $s_f:=\tilde f\circ s$ und $\Lambda_f:U\to G$ durch $s_f=\big(s\cdot\Lambda_f\big)\circ f^{-1}$ bestimmt ist.
\emph{In Minkowski} mit globalem konstantem Rahmen und linearem $f=(a,\Lambda)$ vereinfacht sich
\[
(f^*_\rho\phi)(x)=\rho(\Lambda)\,\phi(\Lambda x + a).
\]

\emph{Differentialoperatoren: Natürlichkeit, Symbol, Verbindungen}

\emph{Kovariante Ableitung und ihr Transformationsgesetz.}
Eine $G$-Verbindung auf $\mathrm{Fr}_\eta(M)$ ist eine $\frak g$-wertige 1-Form $\omega$,
lokal $\omega^{[s]}:=s^*\omega\in\Omega^1(U,\frak g)$.
Sie induziert auf $E_\rho$ die kovariante Ableitung
\[
\nabla^{[s]}\phi^{[s]}=d\phi^{[s]}+\rho_*\!\big(\omega^{[s]}\big)\,\phi^{[s]}.
\]
Unter einem \emph{Rahmenwechsel} $s\mapsto s\cdot\Lambda$ gilt das Eichgesetz
\[
\omega^{[s]}\mapsto \omega^{[s']}=\Lambda^{-1}\omega^{[s]}\Lambda+\Lambda^{-1}d\Lambda,
\qquad
\nabla^{[s']}\phi^{[s']}=\rho(\Lambda)\,\big(\nabla^{[s]}\phi^{[s]}\big).
\]
Unter einer \emph{Isometrie} $f$ gilt (mit $\tilde f^*\omega=\omega$ für Levi-Civita)
\[
f_\rho^*(\nabla\phi)=\nabla(f_\rho^*\phi),
\]
d.\,h. Pullback und kovariante Ableitung \emph{kommutieren}.

\emph{Natürliche Differentialoperatoren.}
Ein lokaler Differentialoperator $D:\Gamma(E_\rho)\to\Gamma(E_\sigma)$ heißt \emph{natürlich} (Lorentz-kovariant), wenn
\[
\boxed{\quad f^*_\sigma\circ D \;=\; D\circ f^*_\rho \quad\ \text{für alle }f\in\Iso(M,g).}
\]
In einem Rahmen $s$ bedeutet dies Kompatibilität mit den Transformationsgesetzen von
$\phi^{[s]}$, $\omega^{[s]}$ und den Metrik-/Clifford-Daten.

\emph{Hauptsymbol.}
Für einen Operator $D$ erster Ordnung ist sein Hauptsymbol $\sigma_1(D):T^*M\otimes E_\rho\to E_\sigma$
bündelmorphismus-linear. Natürlichkeit impliziert
\[
\sigma_1(D)_{f(p)}\big((d f_p^{-1})^T\xi\big)
\;=\; \rho_\sigma(\Lambda_f(p))\,\sigma_1(D)_p(\xi)\,\rho_\rho(\Lambda_f(p))^{-1}.
\]
Beispiel: Dirac-Operator $D_{\text{Dir}}=i\slashed{\nabla}-m$ hat
$\sigma_1(D_{\text{Dir}})(\xi)=\gamma(\xi^\sharp)$ und ist natürlich, weil
$\gamma$ und $\nabla$ aus $g$ und der Spinverbindung stammen.

\emph{Kanonische Beispiele (SRT, global inertial).}
\begin{itemize}
\item \textbf{Skalar (Klein–Gordon):} $\rho=\mathbf{1}$, $E_\rho=M\times\C$,
\[
D_{\mathrm{KG}}=\Box_\eta+m^2,\qquad
\big(D_{\mathrm{KG}}\circ f^*\big)\phi=\big(f^*\circ D_{\mathrm{KG}}\big)\phi.
\]
In Komponenten: $\partial'_\mu=(\Lambda^{-1})^\nu{}_\mu\partial_\nu$, $\Lambda^T\eta\Lambda=\eta$.

\item \textbf{Vektor-/1-Form-Feld (Maxwell in Potentialform):}
$A\in\Omega^1(M)$, $F=dA$, $D_{\mathrm{Max}}=d^\dagger d$.
Isometrien sind $g$-verträglich $\Rightarrow$ $d$ und $d^\dagger$ sind natürlich,
also $f^* D_{\mathrm{Max}}=D_{\mathrm{Max}} f^*$.

\item \textbf{Spinorfeld (Dirac):}
Spin-Struktur $\widetilde{\mathrm{Fr}}_\eta(M)\to\mathrm{Fr}_\eta(M)$,
$\rho:\Spin(1,3)\to\GL(\C^4)$, $\psi\in\Gamma(S)$.
Für $f$ mit Lift $\tilde f_{\Spin}$ und zugehörigem $S(\Lambda_f)$ gilt
\[
\psi\mapsto f^*_\rho\psi,\quad
(i\gamma^\mu\nabla_\mu-m)(f^*_\rho\psi)=f_\rho^*((i\gamma^\mu\nabla_\mu-m)\psi).
\]
In Minkowski und globalem Inertialrahmen: $\psi'(x)=S(\Lambda)\psi(\Lambda x+a)$.
\end{itemize}

\emph{Distributionen, Greensche Operatoren, Impulsraum}

\emph{Distributionelle Schnitte.}
Setze $\Gamma^{-\infty}(E_\rho):=\big(\Gamma_c^\infty(E_\rho^\vee\otimes\mathrm{Vol}_g)\big)^\prime$.
Ist $D$ natürlich und formal adjungiert $\!D^\dagger$, so kommutiert der Pullback auch im distributionellen Sinn:
\[
\langle f^*_\sigma D\phi,\,\varphi\rangle
=\langle \phi,\, D^\dagger f_*^\sigma\varphi\rangle
=\langle \phi,\, f_*^\rho D^\dagger\varphi\rangle
=\langle D\phi,\, f_*^\sigma\varphi\rangle
=\langle D f^*_\rho\phi,\,\varphi\rangle.
\]

\emph{Greensche Operatoren.}
Für hyperbolische, natürliche Operatoren (KG, Dirac) sind die kausalen Greenschen Operatoren
$E^{\mathrm{ret/adv}}$ ebenfalls natürlich:
\[
f^*_\rho\circ E^{\mathrm{ret/adv}} \;=\; E^{\mathrm{ret/adv}}\circ f^*_\rho.
\]

\emph{Impulsraum (Minkowski).}
Für temperierte Felder $\phi$ bewirkt $f=(a,\Lambda)$
\[
\widehat{(f^*_\rho\phi)}(p)=\rho(\Lambda)\,e^{+i\,a\cdot p}\,\hat\phi(\Lambda^{-1}p),
\]
die Massenschale $p^2=m^2$ bleibt invariant; das ist die unitäre (projektive) Darstellung der Poincaré-Gruppe
im Einteilchenbild.

\emph{„Arbeitsrezept“ (Beobachterwechsel im Detail)}

\begin{enumerate}
\item \textbf{Feldtyp fixieren:} Wähle $\rho:G\to\GL(V)$ und das assoziierte Bündel $E_\rho$.
\item \textbf{Schnittsicht:} Schreibe $\phi\in\Gamma(E_\rho)$ als äquivariante $\Phi:\mathrm{Fr}_\eta(M)\to V$.
\item \textbf{Passiver Wechsel (Rahmen):} $s\mapsto s\cdot\Lambda$ induziert
$\phi^{[s]}\mapsto \phi^{[s']}(p)=\rho(\Lambda(p))\,\phi^{[s]}(p)$,
und für jede kovariante Ableitung $\nabla$:
$\nabla^{[s']}\phi^{[s']}=\rho(\Lambda)\,\nabla^{[s]}\phi^{[s]}$.
\item \textbf{Aktiver Wechsel (Isometrie):} $f\in\Iso(M,g)$ mit Lift $\tilde f$ liefert
$f^*_\rho\phi$ via $\Phi\mapsto \Phi\circ\tilde f$.
In Minkowski: $(f^*_\rho\phi)(x)=\rho(\Lambda)\,\phi(\Lambda x+a)$.
\item \textbf{Operatorverhalten:} Für natürliche $D$ gilt stets
\[
\boxed{\,f^*_\sigma\circ D = D\circ f^*_\rho\,}.
\]
Lokal: Rahmenwechsel-/Eichgesetz für $\omega^{[s]}$, Kettenregel für Ableitungen,
und $G$-Kovarianz der Darstellung reichen aus.
\end{enumerate}

\emph{Zwei Mikro-Beispiele}

\emph{(i) Klein–Gordon (Skalar).}
$E_{\mathbf 1}=M\times\C$, $\nabla=d$, $D=\Box_\eta+m^2$.
Für $f=(a,\Lambda)$:
\[
\phi'(x):=(f^* \phi)(x)=\phi(\Lambda x+a),\qquad
(\Box_\eta+m^2)\phi'= \big((\Box_\eta+m^2)\phi\big)\!\circ f.
\]

\emph{(ii) Dirac (Spinor).}
$S\to M$ Spinorbündel, $\nabla^S$ Spinverbindung, $D=i\slashed{\nabla} - m$.
In Minkowski (globaler Vierbein, $\omega=0$):
\[
\psi'(x)=S(\Lambda)\,\psi(\Lambda x+a),\qquad
(i\gamma^\mu\partial_\mu-m)\psi' = S(\Lambda)\,(i\gamma^\mu\partial_\mu-m)\psi\circ f.
\]

\emph{Fazit}

Bezugssystemwechsel sind \emph{entweder} (passiv) Rahmenwechsel im Hauptbündel $\mathrm{Fr}_\eta(M)$
\emph{oder} (aktiv) Isometrien des Basisraums mit induzierter Lift-Wirkung.
Felder sind Schnitte assoziierter Bündel bzw. äquivariante Funktionen auf $\mathrm{Fr}_\eta(M)$.
Natürliche (geometrisch konstruierte) Differentialoperatoren kommutieren
mit beiden Arten des Wechsels. In der SRT (Minkowski, global inertial)
reduziert sich der gesamte Apparat auf die üblichen, komponentenweisen Lorentz- und
Translationsgesetze.

\pagebreak

\subsection{Aktive vs. Passive Transformationen – konzeptionelle Trennung und physikalische Bedeutung}

\emph{1. Grundidee.}
In der formalen bündeltheoretischen Sprache sind \emph{passive} und \emph{aktive} Transformationen
zwei konzeptionell verschiedene, aber mathematisch eng verwandte Vorgänge:

\begin{itemize}
  \item Der \textbf{passive Wechsel} betrifft die \emph{Darstellung desselben Feldes}
  in einem \emph{anderen lokalen Rahmen} oder Koordinatensystem.
  Er verändert nicht das Feld als geometrisches Objekt, sondern nur seine \emph{Komponentendarstellung}.

  \item Der \textbf{aktive Wechsel} betrifft eine \emph{tatsächliche Transformation des Feldes selbst}
  durch eine Isometrie der Raumzeit.
  Er beschreibt die \emph{Abbildung physikalischer Konfigurationen auf neue Konfigurationen}
  – etwa das „Verschieben“ oder „Verkippen“ eines Feldes im Raumzeitkontinuum.
\end{itemize}

Beide Transformationen sind formal durch dieselbe Gruppenstruktur
($SO^+(1,3)$ bzw. Poincaré-Gruppe) beschrieben,
aber sie wirken auf \emph{verschiedenen Ebenen der Beschreibung}.

\bigskip
\emph{2. Passiver Bezugssystemwechsel (Rahmenwechsel).}

Sei $s:U\to \mathrm{Fr}_\eta(M)$ ein lokaler orthonormaler Rahmen.
Das Feld $\phi\in\Gamma(E_\rho)$ erscheint darin als lokale Darstellung
$\phi^{[s]}:U\to V$ mit
\[
\phi(p) = [\,s(p),\,\phi^{[s]}(p)\,].
\]
Wählt man nun einen neuen Rahmen $s' = s\cdot\Lambda$ mit glatter $\Lambda:U\to SO^+(1,3)$, so
ändert sich die lokale Darstellung zu
\[
\boxed{\quad
\phi^{[s']}(p)=\rho\!\big(\Lambda(p)\big)\,\phi^{[s]}(p).
\quad}
\]
Dabei bleibt der geometrische Schnitt $\phi$ \emph{identisch}.
Es wird lediglich eine andere \emph{Komponentenbeschreibung} gewählt.

\smallskip
\emph{Anschaulich:}
Man „dreht oder boostet das Koordinatengitter“, ohne das physikalische Objekt selbst zu verändern.
Wie in der linearen Algebra bleibt der geometrische Vektor derselbe, aber seine Koordinaten ändern sich.

\smallskip
\emph{Formales Verhalten der Ableitung:}
Für eine kovariante Ableitung $\nabla$ gilt das Eichgesetz
\[
\nabla^{[s']}\phi^{[s']}=\rho(\Lambda)\,\nabla^{[s]}\phi^{[s]}.
\]
Das heißt: die Gleichungsform $D(\phi)=0$ bleibt erhalten, wenn Operatoren und Felder
gemeinsam in den neuen Rahmen überführt werden.

\bigskip
\emph{3. Aktiver Bezugssystemwechsel (Isometrie).}

Eine Isometrie $f\in\Iso(M,g)$ beschreibt einen \emph{aktiven Vorgang}:
das Feld selbst wird in der Raumzeit verschoben oder transformiert.
Die induzierte Bündelabbildung $\tilde f:\mathrm{Fr}_\eta(M)\to \mathrm{Fr}_\eta(M)$ lautet
\[
\tilde f(p,e)=(f(p),df_p\circ e).
\]
Der Pullback $f^*_\rho$ wirkt dann auf einem Schnitt $\phi$ (äquivalent auf der äquivarianten Darstellung
$\Phi:\mathrm{Fr}_\eta(M)\to V$) durch
\[
(f^*_\rho\phi)\;\widehat{=}\;\Phi\circ\tilde f,
\quad\text{d.\,h.}\quad
(f^*_\rho\phi)(p)=[\,e,\;\Phi(\tilde f(e))\,].
\]
In einem lokalen Trivialisat (mit Rahmen $s$) liest sich das konkret als
\[
\boxed{\quad
(f^*_\rho\phi)^{[s]}(p)
=\rho\!\big(\Lambda_f(p)\big)\,\phi^{[s_f]}\!\big(f(p)\big),
\quad s_f=\tilde f\circ s.
\quad}
\]

\smallskip
\emph{Physikalisch:}
$f$ beschreibt hier einen realen \emph{Beobachterwechsel im physikalischen Sinne}:
man betrachtet dieselbe physikalische Situation \emph{aus einem anderen Inertialsystem}.
Dabei werden:
\begin{itemize}
  \item die \emph{Ereignispunkte} verschoben: $p\mapsto f(p)$,
  \item das \emph{Feld} entlang dieser Verschiebung mittransportiert (Pullback),
  \item die \emph{internen Komponenten} des Feldes durch $\rho(\Lambda_f)$ transformiert.
\end{itemize}

\smallskip
Im Minkowski-Fall $(M,g)=(\R^{1,3},\eta)$ mit $f(x)=\Lambda x + a$
und konstanter Darstellung $\rho(\Lambda)$ vereinfacht sich das zu:
\[
\boxed{\quad
(f^*_\rho\phi)(x)=\rho(\Lambda)\,\phi(\Lambda x + a).
\quad}
\]
Dies ist genau die übliche \emph{aktive Poincaré-Transformation}
eines Feldes in der Physik.

\bigskip
\emph{4. Vergleich und physikalische Relevanz.}

\begin{center}
\renewcommand{\arraystretch}{1.4}
\begin{tabular}{@{}lll@{}}
\textbf{Aspekt} & \textbf{Passiver Wechsel} & \textbf{Aktiver Wechsel} \\
Objekt & Darstellung $\phi^{[s]}$ & Feld $\phi$ selbst \\
Wirkung & $s\mapsto s\cdot\Lambda$ & $f:M\to M$, $\tilde f$ auf Bündel \\
Verändert & nur Komponenten & Argument und interne Struktur \\
Interpretation & Änderung des Koordinatensystems & reale Transformation des physikalischen Systems \\
Mathematisch & lokale Eichtransformation & Pullback durch Isometrie \\
In Minkowski & \textit{„Wechsel der Messbasis“} & \textit{„Poincaré-Aktion auf Feldern“} \\
\end{tabular}
\end{center}

\smallskip
\emph{Physikalisch relevant:}
In der \textbf{speziellen Relativitätstheorie} (flacher Minkowski-Raum)
interessiert man sich primär für den \emph{aktiven Wechsel}, 
weil er die tatsächliche Symmetrie der Naturgesetze beschreibt:
\[
D(\phi)=0 \quad\Longleftrightarrow\quad D(f^*_\rho\phi)=0,
\quad f\in\Iso(\R^{1,3},\eta).
\]
Der passive Wechsel ist dagegen eine rein \emph{deskriptive Freiheit}:
er ändert nur, wie man dieselbe Situation im selben System beschreibt.

\bigskip
\emph{5. Informationsfluss bei aktiver Transformation.}

Sei $\phi(x)$ ein Feld in einem Inertialsystem $S$.  
Ein zweiter Beobachter $S'$ bewegt sich relativ zu $S$ durch $f(x)=\Lambda x+a$.
Dann gilt:

\[
\begin{aligned}
&\text{(i) Ereigniszuordnung:} &&x' = f^{-1}(x) = \Lambda^{-1}(x - a),\\[2pt]
&\text{(ii) Feldzuordnung:} &&\phi'(x') = \rho(\Lambda)\,\phi(x),\\[2pt]
&\text{(iii) Informationsfluss:} &&\phi'(x')\ \text{trägt dieselbe physikalische Information über Ereignis }x=f(x'). 
\end{aligned}
\]
Somit beschreibt die aktive Transformation die \emph{Übertragung physikalischer Information}
zwischen Beobachtern: jedes Ereignis $x$ im alten System entspricht einem Ereignis $x'$
im neuen, und das Feld $\phi'(x')$ kodiert die gleiche physikalische Realität, nur in
anderen Koordinaten und mit transformierten internen Komponenten.


\pagebreak
\subsection{Lorentztransformationen als Rechts-/Linkswirkungen auf Bündeln}

\emph{Framenbündel und Strukturgruppe}
Sei $(M,g)$ eine Lorentz-Mannigfaltigkeit (in der SRT: $M=\mathbb{R}^{1,3}$, $g=\eta$).
Das \emph{Lorentz-Rahmenbündel}
\[
\mathrm{Fr}_\eta(M)
:=\bigl\{(p,e)\ \big|\ p\in M,\ e=(e_0,\dots,e_3)\text{ Basis von }T_pM,\ g(e_\mu,e_\nu)=\eta_{\mu\nu}\bigr\}
\]
ist ein Hauptbündel über $M$ mit Strukturgruppe $G:=O(1,3)$ (bzw.\ $SO^+(1,3)$ in der orientierten/zeitorientierten Komponente).


\emph{Rechtswirkung der Lorentzgruppe}
Die Gruppenwirkung
\[
R_\Lambda:\ \mathrm{Fr}_\eta(M)\to\mathrm{Fr}_\eta(M),\qquad
(p,e)\mapsto(p,e\cdot\Lambda),\quad (e\cdot\Lambda)_\mu:=e_\nu\,\Lambda^\nu{}_\mu,\ \ \Lambda\in G,
\]
ist frei und transitiv auf jeder Faser und \emph{definiert} die Hauptbündelstruktur. 
Diese \emph{Rechtswirkung} modelliert den \textbf{passiven Rahmenwechsel} (Bezugssystemwechsel am selben Raumzeitpunkt).


\emph{Linkswirkung durch Isometrien der Basis}
Jede Isometrie $f\in\Iso(M,g)$ hebt sich kanonisch zu einer \emph{Linkswirkung}
\[
\tilde f:\ \mathrm{Fr}_\eta(M)\to\mathrm{Fr}_\eta(M),\qquad
\tilde f(p,e):=(f(p),\,df_p\circ e)
\]
an. Es gilt die Kommutativität mit der Rechtswirkung:
\[
\tilde f\big((p,e)\cdot\Lambda\big)=\tilde f(p,e)\cdot\Lambda,\qquad \forall\,\Lambda\in G.
\]
Diese \emph{Linkswirkung} entspricht dem \textbf{aktiven Beobachterwechsel} (Raumzeit-Isometrie auf der Basis).


\emph{Assoziierte Bündel und Felder}
Sei $\rho:G\to\GL(V)$ eine Darstellung. Das \emph{assoziierte Bündel}
\[
E_\rho\ :=\ \mathrm{Fr}_\eta(M)\times_\rho V\ \longrightarrow\ M,\qquad
(e,v)\sim(e\cdot\Lambda,\ \rho(\Lambda)^{-1}v),
\]
trägt natürliche Links-/Rechtswirkungen:
\begin{itemize}
  \item \textbf{Rechts (Lorentz, passiv):} $R_\Lambda:[e,v]\mapsto[e\cdot\Lambda,\ v]$, was auf der Faser den Komponentenwechsel $v\mapsto \rho(\Lambda)\,v$ in einer gewählten lokalen Rahmensektion induziert.
  \item \textbf{Links (Isometrie, aktiv):} $\tilde f:[e,v]\mapsto[\tilde f(e),\ v]$, was dem Pullback $f^*$ von Schnitten entspricht.
\end{itemize}
Ein \emph{Feld} vom Typ $\rho$ ist ein Schnitt $\phi\in\Gamma(E_\rho)$, äquivalent eine äquivariante Abbildung
\[
\Phi:\mathrm{Fr}_\eta(M)\to V,\qquad \Phi(e\cdot\Lambda)=\rho(\Lambda)^{-1}\Phi(e).
\]


\emph{Aktiv vs.\ Passiv auf Feldern}
Für $\phi\in\Gamma(E_\rho)$, dargestellt durch eine äquivariante $\Phi$, gilt:
\begin{enumerate}
  \item \textbf{Passiver Rahmenwechsel (Rechtswirkung).} Wählt man lokal einen Rahmen $s:U\to\mathrm{Fr}_\eta(M)$, so sind die Komponenten $\phi^{[s]}:U\to V$ durch $\phi^{[s]}(p)=\Phi(s(p))$ gegeben und transformieren bei $s\mapsto s\cdot\Lambda$ punktweise durch
  \[
  \phi^{[s\cdot\Lambda]}(p)=\rho\big(\Lambda(p)\big)\,\phi^{[s]}(p).
  \]
  \item \textbf{Aktiver Beobachterwechsel (Linkswirkung).} Für $f\in\Iso(M,g)$ ist das transformierte Feld
  \[
  (f^*\phi)(p)\ :=\ [\tilde f^{-1}(s(p)),\,\phi^{[s]}(f(p))]
  \quad\Longleftrightarrow\quad
  \Phi\ \mapsto\ \Phi\circ\tilde f.
  \]
\end{enumerate}


\emph{Kovarianz natürlicher Differentialoperatoren}
Sei $D:\Gamma(E_\rho)\to\Gamma(E_\sigma)$ ein \emph{natürlicher} Differentialoperator (aus $g$ und der zugehörigen Verbindung gebaut). Dann gilt für alle $f\in\Iso(M,g)$
\[
f^*\circ D\ =\ D\circ f^*.
\]
In lokaler Rahmendarstellung bedeutet dies Rahmenwechsel-Kovarianz: wechselt $s\mapsto s\cdot\Lambda$, so transformieren die Komponenten von $D$ darstellungsgetreu, die \emph{Operatorform} bleibt erhalten.


\emph{Minkowski-Spezialfall}
Für $(M,g)=(\mathbb{R}^{1,3},\eta)$ ist das Rahmenbündel trivial:
\[
\mathrm{Fr}_\eta(\mathbb{R}^{1,3})\cong \mathbb{R}^{1,3}\times SO^+(1,3).
\]
\begin{itemize}
  \item \textbf{Rechts (passiv):} $(x,e)\cdot\Lambda=(x,e\Lambda)$ entspricht dem klassischen Lorentz-\emph{Komponentenwechsel} im selben Ereignis.
  \item \textbf{Links (aktiv, Poincaré):} Für $(a,\Lambda)$ wirkt
  \[
  (a,\Lambda)\cdot(x,e)\ =\ \big(\Lambda x+a,\ \Lambda e\big),
  \]
  und auf einem $\rho$-Feld
  \[
  \phi(x)\ \longmapsto\ \big((a,\Lambda)^*\phi\big)(x)\ =\ \rho(\Lambda)\,\phi(\Lambda^{-1}(x-a)).
  \]
  \item Natürliche Operatoren (z.\,B.\ Klein–Gordon $D=\Box_\eta+m^2$, Dirac $D=i\gamma^\mu\partial_\mu-m$, Maxwell in Potentialform $D=d^\dagger d$) erfüllen
  \[
  (a,\Lambda)^*\circ D\ =\ D\circ (a,\Lambda)^*.
  \]
\end{itemize}


\emph{Konzeptuelle Zuordnung}
\begin{itemize}
  \item \textbf{Rechtswirkung} $R_\Lambda$ auf $\mathrm{Fr}_\eta(M)$: \emph{passiver} Wechsel des lokalen orthonormalen Rahmens (Darstellung/Komponentenwechsel).
  \item \textbf{Linkswirkung} $\tilde f$ durch $f\in\Iso(M,g)$: \emph{aktiver} Beobachterwechsel (Isometrie der Basis, Pullback $f^*$ auf Feldern).
  \item Die \textbf{Forminvarianz} der Gesetze ist die \textbf{Natürlichkeit} der entsprechenden Differentialoperatoren bzgl.\ der Linkswirkung, kompatibel mit der darstellungsinduzierten Rechtswirkung.
\end{itemize}


\pagebreak

\subsection{Aktive Poincaré-Transformationen auf skalaren Feldern im Minkowski-Raum}

\emph{Setting und Notation.}
Sei $(M,g)=(\mathbb{R}^{1,3},\eta)$ mit $\eta=\mathrm{diag}(1,-1,-1,-1)$.
Die Poincaré-Gruppe ist das semidirekte Produkt
\[
\mathcal{P}\;=\;\mathbb{R}^{1,3}\rtimes SO^+(1,3)
\;=\;\{(a,\Lambda)\mid a\in\mathbb{R}^{1,3},\ \Lambda\in SO^+(1,3)\},
\]
mit Gruppenmultiplikation
\[
(a_2,\Lambda_2)\cdot(a_1,\Lambda_1)\;=\;(a_2+\Lambda_2 a_1,\ \Lambda_2\Lambda_1),
\qquad (0,\mathbf{1})\ \text{neutral},\quad (a,\Lambda)^{-1}=(-\Lambda^{-1}a,\Lambda^{-1}).
\]

\emph{Aktive Wirkung auf der Basis (Ereignisraum).}
Die \emph{aktive} Poincaré-Wirkung auf $M$ ist die linke Gruppenwirkung
\[
\begin{aligned}
\Phi:\ \mathcal{P}\times M&\longrightarrow M,\\
\bigl((a,\Lambda),x\bigr)&\longmapsto \Phi\bigl((a,\Lambda),x\bigr):=\Lambda x + a.
\end{aligned}
\]
Damit wird ein Ereignis $x$ in ein \emph{neues} Ereignis $\Lambda x+a$ überführt.

\emph{Skalares Feld als Schnitt.}
Ein (komplexes) skalares Feld ist ein Schnitt des trivialen Linienbündels
$E:=M\times\mathbb{C}\to M$; äquivalent ist $\Gamma(E)\cong C^\infty(M,\mathbb{C})$.
Wir schreiben ein Feld als Funktion $\phi:M\to\mathbb{C}$, $x\mapsto \phi(x)$.

\emph{Induzierte Links\-wirkung auf Schnitten (Pullback-Definition).}
Die aktive Basiswirkung $\Phi$ induziert eine \emph{linke} Gruppenwirkung
\[
\Pi:\ \mathcal{P}\times\Gamma(E)\longrightarrow\Gamma(E),
\]
definiert durch \emph{Pullback entlang der Basiswirkung}:
\[
\boxed{\quad
\bigl(\Pi((a,\Lambda),\phi)\bigr)(x)
:= \phi\!\bigl(\Phi((a,\Lambda),x)\bigr)
= \phi(\Lambda x + a).
\quad}
\]
Kurzschreibweise: $(a,\Lambda)^*\,\phi := \phi\circ(\Lambda(\cdot)+a)$.
Dies ist genau die aktive Poincaré-Transformation eines \emph{skalaren} Feldes.

\emph{Axiome einer Gruppenwirkung (Nachweis).}
Für alle $(a_i,\Lambda_i)\in\mathcal{P}$, $\phi\in\Gamma(E)$ und $x\in M$ gilt:
\[
\begin{aligned}
\text{(i) Einheit:}\quad & (0,\mathbf{1})^*\phi(x)=\phi(\mathbf{1}\,x+0)=\phi(x),\\[2pt]
\text{(ii) Verträglichkeit:}\quad &
\bigl((a_2,\Lambda_2)\cdot(a_1,\Lambda_1)\bigr)^*\phi(x)
= \phi\bigl(\Lambda_2(\Lambda_1 x+a_1)+a_2\bigr)\\
&= \bigl((a_1,\Lambda_1)^*\phi\bigr)\bigl(\Lambda_2 x+a_2\bigr)
= (a_2,\Lambda_2)^*\bigl((a_1,\Lambda_1)^*\phi\bigr)(x).
\end{aligned}
\]
Also ist $\Pi$ eine wohldefinierte \emph{linke} Gruppenwirkung.

\emph{Abbildungsvorschrift (explizit zum Ausrechnen).}
Um $(a,\Lambda)^*\phi$ an einem Punkt $x$ zu berechnen:
\begin{enumerate}
  \item \textbf{Bildpunkt bestimmen:} $x'=\Lambda x + a$.
  \item \textbf{Feld am Bildpunkt auswerten:} $((a,\Lambda)^*\phi)(x)=\phi(x')$.
\end{enumerate}
Für Ableitungen folgt aus der Kettenregel:
\[
\partial_\mu\bigl((a,\Lambda)^*\phi\bigr)(x)
= \Lambda^\nu{}_{\mu}\,(\partial_\nu\phi)(\Lambda x+a),\qquad
\square_\eta\bigl((a,\Lambda)^*\phi\bigr)=\bigl(\square_\eta\phi\bigr)\circ(\Lambda(\cdot)+a).
\]
Insbesondere kommutiert der Klein–Gordon-Operator $P_m=\square_\eta+m^2$ mit der Wirkung:
\[
(a,\Lambda)^*\,P_m\phi \;=\; P_m\bigl((a,\Lambda)^*\phi\bigr).
\]

\emph{Alternativ (gleichwertige „Inertialbeobachter“-Schreibweise).}
Man stellt oft auf „neue“ Koordinaten $x':=\Lambda x + a$ um und definiert das transformierte Feld
\[
\phi'(x')\ :=\ (a,\Lambda)^*\phi\bigl(\Lambda^{-1}(x'-a)\bigr)
=\phi(x')\quad\text{(skalar!)}.
\]
Schreibt man alles wieder als Funktion von $x$, erhält man die übliche kompakte Formel
\[
\boxed{\quad \phi'(x) = \bigl((a,\Lambda)^*\phi\bigr)(x)
\;=\; \phi(\Lambda x + a).\quad}
\]
Beide Notationen sind identisch; sie unterscheiden nur, ob man $x$ oder $x'$ als freies Argument wählt.

\emph{Bemerkung (Darstellungen und Spin).}
Für Nicht-Skalarfelder (Vektoren/Spinoren) tritt zusätzlich die \emph{interne} Darstellung
$\rho(\Lambda)$ in der Faser auf; für Skalarfelder ist $\rho=\mathbf{1}$ trivial und obige Formel
reduziert sich auf den reinen Pullback auf der Basis. Deshalb steht beim Skalarfeld \emph{kein}
präfaktor $\rho(\Lambda)$ vor $\phi$.

\emph{Verallgemeinerung (Distributionen).}
Die gleiche Vorschrift definiert eine stetige linke Gruppenwirkung auf temperierten
Distributionen $\mathcal{S}'(M)$ durch Dualität:
\[
\langle (a,\Lambda)^*u,\; \varphi\rangle
:= \langle u,\; \varphi\circ(\Lambda^{-1}(\cdot - a)) \rangle,
\qquad u\in\mathcal{S}'(M),\ \varphi\in\mathcal{S}(M).
\]
Bei glattem $u=\phi$ fällt dies mit der obigen Formel zusammen.

\emph{Zusammenfassung (kompakte Definition).}
Die \emph{aktive Poincaré-Transformation} eines skalaren Feldes
$\phi\in\Gamma(M\times\mathbb{C})\cong C^\infty(M)$ ist die linke Gruppenwirkung
\[
\boxed{\;
\begin{aligned}
\Pi:\ \mathcal{P}\times C^\infty(M)&\longrightarrow C^\infty(M),\\
\bigl((a,\Lambda),\phi\bigr)&\longmapsto (a,\Lambda)^*\phi:=\phi\circ(\Lambda(\cdot)+a),
\end{aligned}
\;}
\]
die die Gruppenaxiome erfüllt und mit allen natürlichen (Poincaré-kovarianten)
Differentialoperatoren (z.\,B.\ $P_m$) kommutiert.

\pagebreak
\printbibliography
\end{document}